\section{Deskripsi Tugas}
\subsection{Latar Belakang}
Battlecode merupakan permainan strategi waktu nyata di mana pemain menulis kode untuk mengendalikan agen otonom berupa robot. Robot-robot tersebut harus bekerja sama untuk mengalahkan tim lawan dengan mengelola sumber daya, melakukan koordinasi, serta menerapkan strategi yang efektif.

Dalam permainan ini, pemain harus memanfaatkan berbagai konsep ilmu komputer seperti kecerdasan buatan, penelusuran jalur, algoritma terdistribusi, serta mekanisme komunikasi antar agen. Selain itu, setiap robot memiliki batasan sumber daya komputasi pada setiap giliran, sehingga pemain perlu merancang algoritma yang efisien dan mampu mengambil keputusan dengan cepat.

Pada Tugas Besar 1 mata kuliah IF2211 Strategi Algoritma ini, mahasiswa diminta untuk merancang dan mengimplementasikan bot permainan Battlecode 2025 yang memanfaatkan strategi algoritma greedy. Bot yang dibuat nantinya akan dipertandingkan dengan bot dari kelompok lain.

\subsection{Komponen Permainan}

\subsubsection{Objective Permainan}

Tujuan utama permainan Battlecode adalah mewarnai wilayah peta sebanyak mungkin. Tim yang berhasil mewarnai lebih dari 70\% petak yang dapat diwarnai akan memenangkan permainan.

Selain itu, kemenangan juga dapat dicapai dengan menghancurkan seluruh robot dan menara milik tim lawan.

\subsubsection{Aturan Tiebreaker}

Apabila setelah 2000 ronde tidak ada tim yang berhasil mewarnai lebih dari 70\% peta, maka permainan akan berakhir dan pemenang ditentukan berdasarkan kriteria berikut secara berurutan:

\begin{enumerate}
\item Luas area peta yang telah diwarnai
\item Jumlah menara sekutu yang masih hidup
\item Total jumlah uang (chips)
\item Total jumlah cat (paint) yang dimiliki robot dan menara
\item Jumlah robot sekutu yang masih hidup
\item Jika masih seri, pemenang ditentukan secara acak
\end{enumerate}

\subsubsection{Peta Permainan}

Setiap permainan dimainkan pada peta berbentuk grid dua dimensi berbentuk persegi panjang dengan ukuran antara $20 \times 20$ hingga $60 \times 60$ petak.

Koordinat peta dimulai dari titik kiri bawah $(0,0)$ dan meningkat ke arah timur (kanan) dan utara (atas). Setiap lokasi pada peta direpresentasikan oleh objek \texttt{MapLocation} yang menyimpan koordinat $x$ dan $y$.

Untuk menjaga keseimbangan permainan, peta dijamin memiliki sifat simetris, baik melalui rotasi maupun refleksi.

\subsubsection{Jarak Pandang dan Pergerakan}

Robot dapat melihat lingkungan di sekitarnya dalam radius $\sqrt{20}$ atau sekitar 4.47 petak. Robot tidak dapat bergerak melewati menara maupun robot lain.

Peta juga dapat memiliki tembok (walls) yang tidak dapat dilewati maupun diwarnai. Jumlah tembok maksimal adalah 20\% dari seluruh petak pada peta. Selain itu, robot juga tidak dapat melewati atau mewarnai reruntuhan (ruins).

\subsubsection{Tower dan Ruins}

Menara memiliki fungsi sebagai zona spawn tim sekaligus unit pertahanan yang dapat menyerang musuh. Untuk membangun menara, pemain harus mewarnai pola khusus pada area reruntuhan (ruins) berukuran $5 \times 5$.

Setiap permainan dimulai dengan sejumlah reruntuhan yang telah tersedia pada peta. Reruntuhan tidak dapat dipindahkan maupun dihancurkan. Pusat reruntuhan memiliki jarak minimal lima petak satu sama lain, dan area di sekitarnya dijamin bebas dari tembok.

Setiap tim memulai permainan dengan dua menara, yaitu:

\begin{itemize}
\item Menara penghasil uang (Money Tower)
\item Menara penghasil cat (Paint Tower)
\end{itemize}

\subsubsection{Sumber Daya}

Dalam permainan Battlecode terdapat dua sumber daya utama, yaitu \textbf{paint} dan \textbf{chips}.

\begin{itemize}
\item Paint

Paint merupakan sumber daya utama yang digunakan untuk mewarnai wilayah peta dan menentukan skor tim. Setiap robot dan menara memiliki kapasitas penyimpanan paint yang berbeda.Robot yang memiliki paint rendah akan mengalami peningkatan cooldown aksi, bahkan tidak dapat melakukan aksi apabila paint habis sepenuhnya.

\item Chips

Chips digunakan untuk membangun robot maupun menara. Chips dihasilkan secara pasif oleh menara penghasil uang setiap ronde.Setiap tim memulai permainan dengan 2500 chips.

\end{itemize}

\subsubsection{Robot}

Robot merupakan unit utama dalam permainan. Setiap robot dapat melakukan berbagai aksi seperti bergerak, menyerang, melakukan sensing, serta berkomunikasi dengan robot lain.

Setiap robot memiliki batas komputasi yang disebut \textit{bytecode}. Jika robot melebihi batas bytecode pada satu giliran, eksekusi program akan dihentikan sementara dan dilanjutkan pada ronde berikutnya.

Robot memiliki atribut utama berupa:

\begin{itemize}
\item Hitpoints (HP)
\item Kapasitas paint
\item Cooldown aksi dan pergerakan
\item Kemampuan khusus sesuai tipe robot
\end{itemize}

\subsubsection{Jenis Robot}

Permainan Battlecode menyediakan beberapa jenis robot dengan karakteristik berbeda.

\begin {itemize}
\item Soldier

Robot tipe Soldier memiliki kemampuan menyerang satu petak dan dapat memberikan damage terhadap menara musuh.

\item Splasher

Robot tipe Splasher memiliki serangan area (area of effect) yang mampu mewarnai atau menimpa cat pada beberapa petak sekaligus.

\item Mopper

Robot tipe Mopper memiliki kemampuan untuk menghapus cat milik musuh serta mentransfer paint kepada robot atau menara sekutu.

\end{itemize}

\subsubsection{Komunikasi}

Robot tidak dapat berbagi variabel secara langsung karena setiap robot menjalankan salinan program yang terpisah. Oleh karena itu, komunikasi dilakukan melalui dua mekanisme utama:

\begin{itemize}
\item \textbf{Pesan (messages)} yang dapat dikirim antar robot atau menara
\item \textbf{Penanda (markers)} yang ditempatkan pada peta untuk memberikan informasi kepada robot sekutu
\end{itemize}

Setiap pesan memiliki ukuran maksimal 4 byte dan hanya disimpan selama lima ronde dalam buffer penerima.

\subsubsection{Batasan Program}

Program robot memiliki beberapa batasan yang harus dipatuhi.

\begin{itemize}

\item Bytecode

Setiap robot memiliki batas komputasi sebesar 17500 bytecode per ronde, sedangkan menara memiliki batas 20000 bytecode per ronde.

\item Memori

Penggunaan memori robot tidak boleh melebihi 8 MB selama pertandingan.

\item Execution Time 

Jika total waktu eksekusi sebuah tim melebihi batas yang ditentukan oleh sistem, maka tim tersebut akan otomatis kalah.

\end{itemize}

\subsubsection{Tujuan Tugas Besar}

Pada tugas besar ini, setiap kelompok diminta untuk:

\begin{itemize}
\item Mengembangkan tiga bot Battlecode menggunakan bahasa Java
\item Mengimplementasikan strategi algoritma greedy pada masing-masing bot
\item Menggunakan heuristic yang berbeda untuk setiap strategi greedy
\item Mengoptimalkan strategi agar mampu memenangkan pertandingan
\end{itemize}

Bot yang dikembangkan akan dipertandingkan dengan bot dari kelompok lain dalam kompetisi yang diselenggarakan pada akhir tugas besar.